% PAGES: 88-91
% Continues subsection 2.7.1 "Solving linear systems corresponding to the
% same matrix" (itself a subsection of section 2.7 "Linear solvers using LU
% decomposition with row pivoting"), begun before page 88 (PDF index) in a
% sibling agent's file. Do not open that file; this file starts fresh
% mid-subsection with Table 2.12, which appears complete and self-contained
% at the top of page 88.

\begin{table}[H]
\centering
\caption{Solution of multiple linear systems corresponding to the same matrix}
\fbox{\begin{minipage}{0.8\textwidth}
\textbf{Input:}\\
$A = $ nonsingular square matrix of size $n$ with LU decomposition\\
$b_i = $ column vectors of size $n$, $i=1:p$

\medskip
\textbf{Output:}\\
$x_i = $ solution to $Ax_i=b_i$, $i=1:p$

\medskip
$[P,L,U] = \text{lu\_row\_pivoting}(A)$;\\
for $i = 1:p$\\
\hspace*{1.5em}$y = \text{forward\_subst}(L,Pb_i)$;\\
\hspace*{1.5em}$x_i = \text{backward\_subst}(U,y)$;\\
end
\end{minipage}}
\end{table}

Thus, the operation count for solving the $p$ linear systems using the method from
Table 2.12 is
\[
\frac23n^3+O(n^2) \;+\; p\bigl(n^2+O(n)+n^2+O(n)\bigr) \;=\;
\frac23n^3+2pn^2+O(n^2)+pO(n),
\]
which is smaller than $\frac23pn^3+pO(n^2)$, the operation count required by solving
the $p$ linear systems sequentially; see (2.69).

As an example of solving multiple systems corresponding to the same matrix, we show
how to efficiently compute the inverse of a nonsingular matrix.

Finding the inverse of a nonsingular matrix one entry at a time requires computing
the determinant of the matrix and one minor for every entry of the matrix, for a
total cost of more than $n^2\cdot n!$ operations. This is beyond regular computer
power limits even for small values of $n$, since, from Stirling's formula, $n!$ is
of the order $\frac{n^{n+1/2}}{e^n}$; for example, computing the inverse of a
nonsingular $27\times27$ matrix using this method would require more than $10^{28}$
operations.

Nonetheless, the inverse of an $n\times n$ nonsingular matrix having an LU
decomposition can be computed very efficiently in $\frac83n^3+O(n^2)$ operations as
follows:

Let $A$ be an $n\times n$ nonsingular matrix, and let $A^{-1} = \text{col}(c_k)_{k=1:n}$
be the column form of the inverse matrix of $A$. Then, $AA^{-1}=I$, where
$I=\text{col}(e_k)_{k=1:n}$ is the identity matrix of size $n$. Using (1.11), we find
that $AA^{-1}=\text{col}(Ac_k)_{k=1:n}$, and therefore $AA^{-1}=I$ is equivalent to
\[
Ac_k \;=\; e_k, \quad \forall k=1:n,
\]
which can be solved using the method from Table 2.12 as follows:

\begin{center}
\fbox{\begin{minipage}{0.6\textwidth}
$[P,L,U] = \text{lu\_row\_pivoting}(A)$\\
for $k = 1:n$\\
\hspace*{1.5em}$y = \text{forward\_subst}(L,Pe_k)$\\
\hspace*{1.5em}$c_k = \text{backward\_subst}(U,y)$\\
end\\
$A^{-1} = \text{col}(c_k)_{k=1:n}$
\end{minipage}}
\end{center}

The operation count for computing the inverse matrix using this method is
\begin{equation}
\frac23n^3+O(n^2) \;+\; n\bigl(n^2+O(n)+n^2+O(n)\bigr) \;=\; \frac83n^3+O(n^2);
\end{equation}
see (10.88) in Section 10.2.3 for a proof of the last equality.

\subsection{Finding discount factors using the LU decomposition}

If the number of cash flow dates is equal to the number of bonds (and if there are
no redundancies), then the discount factors corresponding to the cash flow dates can
be uniquely determined from the bond prices by solving a linear system using LU
decomposition with row pivoting; see also section 2.2.1 for an example when forward
substitution is used to determine the discount factors.

For the example below, recall that a semiannual coupon bond with face value \$100,
coupon rate $C$, and maturity $T$ pays the holder of the bond a coupon payment equal
to $\frac{C}{2}\cdot100$ every six months, except at maturity. The final payment at
maturity is equal to the face value of the bond plus one coupon payment, i.e.,
$100+\frac{C}{2}100$. For example, a semiannual coupon bond with 15 months to
maturity and coupon rate 5\%, i.e., $C=0.05$, has three cash flow dates, in 3, 9,
and 15 months, corresponding to $t_1=\frac{3}{12}=\frac14$, $t_2=\frac9{12}=\frac34$,
and $t_3=\frac{15}{12}=\frac54$, respectively. The corresponding cash flows are
$c_1=2.5$, $c_2=2.5$, and $c_3=102.5$.

\begin{examplex}
The values of the following coupon bonds with face value \$100 are given:

\begin{center}
\begin{tabular}{ccc}
Bond Type & Coupon Rate & Bond Price \\[0.4em]
11 months semiannual & 4\% & \$101.5 \\
17 months annual & 5\% & \$105.5 \\
23 months semiannual & 3\% & \$101 \\
23 months annual & 6\% & \$106.75 \\
\end{tabular}
\end{center}

The cash flows and the cash flow dates of the bonds are recorded in the table
below:

\begin{center}
\begin{tabular}{|l|l|}
\hline
Bond Type & Cash Flow \& Date \\
\hline
11 months semiannual & \begin{tabular}[t]{@{}l@{}}\$2 in 5 months\\ \$102 in 11
months\end{tabular} \\
\hline
17 months annual & \begin{tabular}[t]{@{}l@{}}\$5 in 5 months\\ \$105 in 17
months\end{tabular} \\
\hline
23 months semiannual & \begin{tabular}[t]{@{}l@{}}\$1.5 in 5 months\\ \$1.5 in 11
months\\ \$1.5 in 17 months\\ \$101.5 in 23 months\end{tabular} \\
\hline
23 months annual & \begin{tabular}[t]{@{}l@{}}\$6 in 11 months\\ \$106 in 23
months\end{tabular} \\
\hline
\end{tabular}
\end{center}

Note that the four bonds have exactly four cash flow dates, i.e., 5 months, 11
months, 17 months, and 23 months, and let $d_1,d_2,d_3,d_4$ be the discount factors
corresponding to these dates. Using the formula (2.11) for bond valuation, we find
that
\begin{align*}
101.5 &= 2d_1+102d_2; \\
105.5 &= 5d_1+105d_3; \\
101 &= 1.5d_1+1.5d_2+1.5d_3+101.5d_4; \\
106.75 &= 6d_2+106d_4.
\end{align*}

This linear system can be written in matrix notation as $Ax=b$, where
\[
A = \begin{pmatrix} 2 & 102 & 0 & 0\\ 5 & 0 & 105 & 0\\ 1.5 & 1.5 & 1.5 & 101.5\\
0 & 6 & 0 & 106\end{pmatrix}; \quad
x = \begin{pmatrix} d_1\\ d_2\\ d_3\\ d_4\end{pmatrix}; \quad
b = \begin{pmatrix} 101.5\\ 105.5\\ 101\\ 106.75\end{pmatrix}.
\]

The solution of $Ax=b$ is obtained using the LU decomposition with row pivoting
linear solver\footnote{We provide more numerical details here. The permutation
matrix $P$, lower triangular matrix $L$, and upper triangular matrix $U$ from the
LU decomposition with row pivoting of $A$ are
\[
P = \begin{pmatrix} 0&1&0&0\\ 1&0&0&0\\ 0&0&1&0\\ 0&0&0&1\end{pmatrix}; \quad
L = \begin{pmatrix} 1&0&0&0\\ 0.4&1&0&0\\ 0.3&0.0147&1&0\\
0&0.0588&-0.0841&1\end{pmatrix};
\]
\[
U = \begin{pmatrix} 5&0&105&0\\ 0&102&-42&0\\ 0&0&-29.3824&101.5\\
0&0&0&114.5345\end{pmatrix}.
\]
Furthermore,
\[
y \;=\; \text{forward\_subst}(L,Pb) \;=\; \begin{pmatrix} 105.5\\ 59.3\\ 68.4779\\
109.0197\end{pmatrix}.
\]} from Table 2.11, i.e.,
\[
x \;=\; \text{linear\_solve\_lu\_row\_pivoting}(A,b) \;=\; \begin{pmatrix} 0.9916\\
0.9757\\ 0.9575\\ 0.9518\end{pmatrix}.
\]

Recall that the discount factors corresponding to the cash flow dates are equal to
the entries of $x$. Thus,
\[
d_1 = \text{Disc}\left(\frac5{12}\right) = 0.9916; \quad
d_2 = \text{Disc}\left(\frac{11}{12}\right) = 0.9757;
\]
\[
d_3 = \text{Disc}\left(\frac{17}{12}\right) = 0.9575; \quad
d_4 = \text{Disc}\left(\frac{23}{12}\right) = 0.9518.
\]

Then, formula (2.12), i.e.,
\[
\text{Disc}(t_i) \;=\; \exp\bigl(-t_ir(0,t_i)\bigr),
\]
can be used to find the corresponding continuously compounded zero rates, as
follows:
\begin{align*}
d_1 = \exp\left(-\frac5{12}r\left(0,\frac5{12}\right)\right) \quad &\Longleftrightarrow
\quad r\left(0,\frac5{12}\right) = -\frac{\ln(d_1)}{5/12} = 0.0202 = 2.02\%; \\
d_2 = \exp\left(-\frac{11}{12}r\left(0,\frac{11}{12}\right)\right) \quad
&\Longleftrightarrow \quad r\left(0,\frac{11}{12}\right) = -\frac{\ln(d_2)}{11/12} =
0.0269 = 2.69\%; \\
d_3 = \exp\left(-\frac{17}{12}r\left(0,\frac{17}{12}\right)\right) \quad
&\Longleftrightarrow \quad r\left(0,\frac{17}{12}\right) = -\frac{\ln(d_3)}{17/12} =
0.0306 = 3.06\%; \\
d_4 = \exp\left(-\frac{23}{12}r\left(0,\frac{23}{12}\right)\right) \quad
&\Longleftrightarrow \quad r\left(0,\frac{23}{12}\right) = -\frac{\ln(d_4)}{23/12} =
0.0257 = 2.57\%.
\end{align*}
\end{examplex}

\section{Cubic spline interpolation}

Interpolating a continuous function from known values at a finite number of
discrete nodes is a frequently used tool for financial applications, e.g., for
inferring a zero--rate curve from known discount factors.

The continuous interpolation problem can be described in general terms as follows:

Find a continuous function $f:[x_0,x_n]\to\R$ such that $f(x_i)=v_i$, for all
$i=0:n$, where $x_0<x_1<\ldots<x_n$, and $v_i$ are the known values of $f(x)$ at
the point $x_i$, for all $i=0:n$.

The simplest way to obtain such a function is by linear interpolation between every
two points $x_{i-1}$ and $x_i$, for $i=1:n$, resulting in the following function:
\[
f(x) \;=\; \frac{(x-x_{i-1})v_i}{x_i-x_{i-1}} + \frac{(x_i-x)v_{i-1}}{x_i-x_{i-1}},
\quad \forall x_{i-1}\le x\le x_i, \quad \forall i=1:n.
\]

However, linear interpolation has a major drawback: the resulting function $f(x)$
is not differentiable at any point $x_i$, $i=1:(n-1)$.

The interpolation method most often used in practice is \textit{cubic spline
interpolation}, which requires the function $f(x)$ to be a polynomial of degree
three (i.e., a cubic polynomial) on each interval $[x_{i-1},x_i]$, for $i=1:n$, as
well as to be twice differentiable on the entire interval $[x_0,x_n]$.

In other words, we are looking for a function $f(x)$ of the form
\begin{equation}
f(x) \;=\; f_i(x) \;=\; a_i+b_ix+c_ix^2+d_ix^3, \quad \forall x_{i-1}\le x\le x_i,
\quad \forall i=1:n,
\end{equation}
such that
\begin{align}
f_i(x_{i-1}) &= v_{i-1}, \quad \forall i=1:n; \\
f_i(x_i) &= v_i, \quad \forall i=1:n; \\
f_i'(x_i) &= f_{i+1}'(x_i), \quad \forall i=1:(n-1); \\
f_i''(x_i) &= f_{i+1}''(x_i), \quad \forall i=1:(n-1).
\end{align}

Using (2.71), we can rewrite (2.72--2.75), respectively, as follows:
\begin{align}
a_i+b_ix_{i-1}+c_ix_{i-1}^2+d_ix_{i-1}^3 &= v_{i-1}, \quad \forall i=1:n; \\
a_i+b_ix_i+c_ix_i^2+d_ix_i^3 &= v_i, \quad \forall i=1:n; \\
b_i+2c_ix_i+3d_ix_i^2 &= b_{i+1}+2c_{i+1}x_i+3d_{i+1}x_i^2, \quad \forall
i=1:(n-1); \\
2c_i+6d_ix_i &= 2c_{i+1}+6d_{i+1}x_i, \quad \forall i=1:(n-1).
\end{align}
